\clearpage
\fancyhead[L]{\small BODE. ASTRONOMISCHES JAHRBUCH.}
\hypertarget{native-568-lower}{}
\begin{center}\small Göttingische gelehrte Anzeigen. Stück 40. Seiten 393..400. 10. März 1814.\\[3pt]\rule{.72\linewidth}{.3pt}\end{center}

\emph{Astronomisches Jahrbuch für das Jahr} 1816, \emph{mit einer Sammlung der neuesten Abhandlungen, Beobachtungen und Nachrichten aus den astronomischen Wissenschaften. Berechnet und herausgegeben von} J.\,E. \textsc{Bode}, königlichem Astronomen und Mitglied der Akademie. 1813. 268 Seiten im Oktavformat.

Das Jahr 1816 zeichnet sich durch eine große Sonnenfinsternis aus (19. November), die in einem Teil Europas für einige Zeit total sein wird, zum Beispiel in Danzig und Warschau. Der Kalender ist in diesem Jahr genauso aufgebaut wie im vorhergehenden. Dem Sonnendurchmesser und der Schiefe der Ekliptik liegen jetzt die neuesten Bestimmungen zugrunde. Bei den Finsternissen der Jupitermonde, die noch nach den Tafeln von \textsc{Wargentin} berechnet sind, ist im Allgemeinen angegeben, wie sich die Unterschiede zwischen diesen Tafeln und denen von \textsc{Delambre} ändern. Die Zahl und der Reichtum der \emph{Zusätze}, die auch diesem Band beigegeben sind, haben uns umso angenehmer überrascht, als wir angesichts der Zeitumstände kaum gehofft hatten, das Jahrbuch überhaupt zur gewohnten Zeit erscheinen zu sehen. Den Anfang macht ein chronologisches Verzeichnis der berühmtesten (verstorbenen) Astronomen seit dem dreizehnten Jahrhundert, ihrer Verdienste, Schriften und Entdeckungen. Darin wird man kaum einen Namen vermissen, der in der Geschichte der Astronomie auch nur irgendwie von Bedeutung ist. — Astronomische Beobachtungen aus den Jahren 1811 und 1812, angestellt an der königlichen Sternwarte in Kopenhagen von Herrn Staatsrat und Ritter \textsc{Bugge}. Oppositionen des Uranus 1811 und 1812, des Saturn 1811 und 1812, der Ceres 1811, der Vesta 1812 und des Mars 1811, Beobachtungen des großen Kometen von 1811, Sternbedeckungen und Finsternisse der Jupitermonde. — Finsternisse der Jupitermonde und Sternbedeckungen, beobachtet an der Sternwarte in Greenwich in den Jahren 1809 und 1810 (aus der gedruckten Sammlung der Beobachtungen \textsc{Maskelyne}s). — Beobachtungen des veränderlichen Sterns $\eta$ im Antinous und Tafeln zur Berechnung seiner größten Helligkeit von Herrn Professor \textsc{Wurm}, eine Fortsetzung des Aufsatzes im Jahrbuch für 1814. Wir wünschen, dass dieser verdiente Astronom uns bald auch mit einer ähnlichen Arbeit über den veränderlichen Stern im Walfisch beschenkt. — Astronomische Beobachtungen aus dem Jahr 1812, angestellt an der königlichen Sternwarte in Berlin. — Beobachtungen der Jupitermonde. — Finsternisse, Sternbedeckungen, Oppositionen der Ceres, Pallas, Juno, Vesta und des Kometen von 1812, beobachtet an der Wiener Sternwarte von Herrn \textsc{Triesnecker}. Meridianbeobachtungen der Ceres und Vesta, des Mars, Saturn und Jupiter aus dem Jahr 1811 von Herrn \textsc{Groombridge}. — Über die Bestimmung der Teilungsfehler eines Spiegelsextanten von Herrn Dr. \textsc{Benzenberg}. Dass man sich auf die absoluten Größen der mit diesen Instrumenten gemessenen Winkel nicht uneingeschränkt verlassen darf, ist wohlbegründet. Dabei kann in ein\SourceLeaf{569}zelnen Fällen offenbleiben, wie viel von dieser Unsicherheit auf wirkliche Teilungsfehler, auf eine geringe Exzentrizität oder darauf zurückgeht, dass die beiden Flächen des großen Spiegels geringfügig von der parallelen Lage abweichen. Das von Herrn B. vorgeschlagene Verfahren, die Fehler des Instruments in einer Tafel zusammenzustellen, lässt sich allerdings nur in sehr ebenen Gegenden anwenden, sofern man nicht zugleich ein Instrument zum Messen kleiner Höhenwinkel hat. Es ist wohl auch nicht überflüssig, dabei zu besonderer Vorsicht zu raten, damit alle Winkel genau von ein und demselben Punkt aus gemessen oder wenigstens auf diesen Punkt reduziert werden. — Verschiedene astronomische Beobachtungen, mitgeteilt von Herrn \textsc{Lee} aus London. Beobachtungen des großen Kometen von 1811 im Ostindischen Ozean an Bord eines Schiffes vom 17. Mai bis zum 15. Juni, ferner verschiedene astronomische Beobachtungen aus Port Jackson in Neuholland und aus Kalkutta. — Astronomische Beobachtungen aus dem Jahr 1812 an der Prager Sternwarte von den Herren \textsc{David} und \textsc{Bittner}. — Beobachtete Zenitabstände der Sonne und Sternbedeckungen in St. Gallen in der Schweiz von Herrn \textsc{von Scherer}. Herr \textsc{Triesnecker} leitet daraus zugleich die geografische Lage dieses Orts ab. — Sichtbare Helligkeitsänderungen des Algol in den Jahren 1814, 1815 und 1816, vorausberechnet von Herrn \textsc{Wurm}. — Beobachtungen des großen Kometen von 1811 vom 6. September bis zum 19. Oktober, Oppositionen des Saturn 1811 und Sternbedeckungen von Herrn \textsc{Derflinger} in Kremsmünster. — Neue Refraktionstafel, aus \textsc{Bradley}s Beobachtungen der Zirkumpolarsterne abgeleitet, von Herrn Professor \textsc{Bessel}. Diese Tafel ist einem Aufsatz des Verfassers im vierten Heft des Königsberger Archivs entnommen. Eine Anleitung zum Gebrauch fehlt, obwohl sie umso nötiger wäre, als das Rechenbeispiel nicht ganz im Sinne des Verfassers ausgeführt ist. — Beobachtungen des großen Kometen im Jahr 1811 mit Bemerkungen über den Aufbau seiner verschiedenen Teile von Herrn Dr. \textsc{Herschel}. Danach folgen Beobachtungen des zweiten Kometen von 1811 mit Bemerkungen über seinen Aufbau vom selben Verfasser. Es handelt sich um Auszüge aus zwei Aufsätzen in den neuesten Bänden der \emph{Philosophical Transactions}. Auf eine ausführliche Besprechung können wir hier umso eher verzichten, als die Originalaufsätze in diesen Blättern bereits früher von einem anderen Rezensenten besprochen worden sind. Wir bemerken daher nur, dass es uns noch viel zu früh scheint, Ergebnisse über die Naturgeschichte dieser bemerkenswerten Himmelskörper festzustellen, so wertvoll \textsc{Herschel}s Beobachtungen dieser Kometen auch sind. Es liegen noch viel zu wenige Beobachtungen vor, und selbst diese sind bei Weitem weniger zuverlässig als astronomische Beobachtungen anderer Art. Nach unserer Überzeugung werden erst die künftigen Jahrhunderte, die die jetzt beobachteten Kometen wiederkehren sehen werden, über das entscheiden können, was jetzt nur eine vage Hypothese sein kann. Herr \textsc{Herschel} stellt unter anderem die Vermutung auf, der Komet von 1807 habe sich, bevor er diesmal sein Perihel erreichte, früher einmal einem anderen Fixstern genähert und sei dadurch gleichsam früher zur Reife gelangt als der von 1811. Er glaubt, die Möglichkeit, dass Kometen andere Sonnen als die unsere umlaufen, werde dadurch sehr wahrscheinlich, dass wir von der großen Zahl der bisher beobachteten erst bei einem einzigen die Wiederkehr mit Gewissheit kennen. Dabei übersieht er aber offenbar, dass nach allen bisherigen Ergebnissen die Umlaufzeiten der Kometen in der Regel eher noch in Jahrtausenden als in Jahrhunderten gemessen werden müssen, dass Kometen erst seit ein paar Jahrhunderten ordentlich beobachtet und ihre Bahnen erst seit wenig mehr als einem Jahrhundert berechnet werden. Was insbesondere den Kometen von 1807 betrifft, so weiß man, dass seine größte Entfernung von der Sonne, ebenso wie die größte Entfernung aller anderen, deren Bahn sich als Ellipse berechnen ließ, gegenüber der Entfernung der nächsten Fixsterne kaum mehr als nichts ist. — Sternbedeckungen, Finsternisse der Jupitermonde, eine Mondfinsternis, Vesta, Uranus und Mars, beobachtet an der Sterwarte in Wilna in den Jahren 1811, 1812 und 1813 von Herrn \textsc{Joh. Sniadecky}. — Beobachtungen des großen Kometen von 1811 an der Sternwarte in Palermo und Berech\SourceLeaf{570}nung seiner Bahnelemente durch Herrn \textsc{Piazzi}, aus dessen gedrucktem Aufsatz. — Aus einem Schreiben von Herrn Professor \textsc{David} in Prag. Es betrifft die abweichenden Ergebnisse, die die Herren \textsc{Vent}, \textsc{Bode} und \textsc{David} zu verschiedenen Zeiten für die Polhöhe der Böhmischen Riesenkuppe gefunden hatten. Darüber steht uns hier natürlich kein Urteil zu. — Nachtrag von Herrn \textsc{von Wisniewsky} zu dem im Jahrbuch für 1813 enthaltenen Verzeichnis der im Jahr 1813 sichtbaren Sternbedeckungen. — Verzeichnis der Länge und Breite von neun der bedeutendsten Fixsterne nach den neuesten Beobachtungen, entnommen aus dem \emph{Nautical Almanac} für 1815. — Geografische Lage verschiedener Orte in Mittelamerika, bestimmt von Herrn \textsc{von Humboldt} (aus seinem \emph{Recueil d'observations astronomiques etc.}). — Projektionsmethode für eine allgemeine Himmelskarte. Sie besteht darin, die Abstände vom Nordpol in der Projektion proportional zu den Tangenten von $\frac{2}{5}$ der Abstände auf der Kugel anzusetzen, während sie bei der stereografischen Projektion bekanntlich proportional zu den Tangenten der halben Abstände auf der Kugel sind. Nach diesem neuen Verfahren hat der Verfasser die Himmelskarte für die neueste Ausgabe seiner geschätzten Anleitung zur Kenntnis des gestirnten Himmels entworfen, deren früheren Ausgaben eine ähnliche Karte in stereografischer Projektion beigegeben war. Er wählte diese Änderung, damit die Grade zum Südpol hin gegenüber den nördlichen weniger unverhältnismäßig anwachsen. Dadurch gehen freilich manche Vorteile der stereografischen Entwurfsart verloren. Das Bild ist dem Original \emph{in den kleinsten Teilen} nicht mehr ähnlich. Alle schief zum Äquator liegenden Kreise werden nicht mehr durch Kreise dargestellt, sondern durch Kurven höherer Ordnung, die sich nur mühsam Punkt für Punkt konstruieren lassen. Die Winkel, unter denen sie sich schneiden, sind nicht mehr gleich denen auf der Kugel. — Geozentrischer Lauf der Pallas vom 1. August 1814 bis zum 1. Februar 1815 und der Vesta vom 18. Oktober 1813 bis zum 9. Juli 1814. — Planetenbeobachtungen an der Sternwarte in Greenwich 1809 und 1810 von Dr. \textsc{Maskelyne}. — Beobachtungen der beiden im Jahr 1813 erschienenen Kometen. Hier sind nur vier Pariser Beobachtungen des ersten Kometen aufgeführt. Über den anderen, den Herr Professor \textsc{Harding} an der hiesigen Sternwarte und zur selben Zeit Herr \textsc{Pons} in Marseille entdeckten, werden nur einige allgemeine Nachrichten gegeben. — Die elliptischen Elemente der Planetenbahnen, für die älteren Planeten aus der dritten Ausgabe der \textsc{Laplace}\emph{'schen exposition du système du monde} (von der 1812 eine vierte Ausgabe erschienen ist), für die neueren nach den neuesten (dem Herausgeber damals bekannten) Bestimmungen von \textsc{Gauss}. Dieser Artikel könnte jetzt schon eine fast vollständige Überarbeitung vertragen, da die Theorien der meisten Planeten seither neu bearbeitet worden sind. — Weitere Beobachtungen des großen Kometen von 1811 mit astronomischen Bemerkungen von Herrn Professor \textsc{Bessel}, eine neue Reduktion der wichtigen Beobachtungen \textsc{von Zach}s zu diesem Kometen im ersten Zweig seiner Bahn. — Beobachtungen des Kometen von 1812 mit den Elementen seiner Bahn von \textsc{Bouvard}. — Methode zur Bestimmung der Abweichung eines Durchgangsinstruments vom Meridian anhand von Beobachtungen der oberen und unteren Kulminationen zweier nördlicher Sterne mit nahezu entgegengesetzter Rektaszension, mit einem Verzeichnis solcher Sterne. Wir halten diese Methode für die zweckmäßigste überhaupt, besonders wenn der eine Stern der Polarstern selbst ist. Wir bemerken nur, dass in diesem Fall der andere Stern ohne Nachteil auch eine nahezu \emph{gleiche} Rektaszension haben darf. — Planetenbeobachtungen an der Pariser Sternwarte im Jahr 1809 von Herrn \textsc{Bouvard} (aus der \emph{Connoissance des tems} für 1812). — Unter den verschiedenen kurzen astronomischen Notizen am Schluss nennen wir hier nur diejenigen, welche die vortrefflichen aus der \textsc{Reichenbach}'sche Werkstatt in München hervorgegangenen astronomischen Instrumenten betreffen, darunter ein 24-füßiger Achromat mit 8 Zoll Öffnung. — Zu diesem Band sind noch zwei Nachträge zu erwähnen. Der erste enthält die Originalbeobachtungen des großen Kometen von 1811, die \textsc{Wisniewsky} im August 1812 angestellt hat, der andere das Fixsternverzeichnis \textsc{Pond}s (wenn auch nicht \SourceLeaf{571}das allerneueste). — Wir können nicht umhin, dem verdienstvollen Herausgeber für den unermüdlichen Eifer zu danken, mit dem er so vielfältiges, für den Astronomen wichtiges Material zusammengetragen hat. Auch dasjenige, was anderen Schriften entnommen ist, die der bei Weitem größten Zahl der Leser jetzt ohnehin meist nicht zugänglich sind, kann nur willkommen sein.
